\documentclass[10pt,a4paper]{amsart}
\usepackage{amsmath,amssymb,amsthm,graphicx}
\usepackage{verbatim}
\newenvironment{reference}{\comment}{\endcomment}
%\newenvironment{reference}{}{}
\newenvironment{slogan}{\comment}{\endcomment}
\newenvironment{history}{\comment}{\endcomment}

% For commutative diagrams we use Xy-pic
\usepackage[all]{xy}

% We use 2cell for 2-commutative diagrams.
\xyoption{2cell}
% Theorem environments.
%
\theoremstyle{plain}
\newtheorem{theorem}[subsection]{Theorem}
\newtheorem{proposition}[subsection]{Proposition}
\newtheorem{lemma}[subsection]{Lemma}

\theoremstyle{definition}
\newtheorem{definition}[subsection]{Definition}
\newtheorem{example}[subsection]{Example}
\newtheorem{exercise}[subsection]{Exercise}
\newtheorem{situation}[subsection]{Situation}

\theoremstyle{remark}
\newtheorem{remark}[subsection]{Remark}
\newtheorem{remarks}[subsection]{Remarks}

\numberwithin{equation}{subsection}

% Macros
%
\def\lim{\mathop{\mathrm{lim}}\nolimits}
\def\colim{\mathop{\mathrm{colim}}\nolimits}
\def\Spec{\mathop{\mathrm{Spec}}}
\def\Hom{\mathop{\mathrm{Hom}}\nolimits}
\def\Ext{\mathop{\mathrm{Ext}}\nolimits}
\def\SheafHom{\mathop{\mathcal{H}\!\mathit{om}}\nolimits}
\def\SheafExt{\mathop{\mathcal{E}\!\mathit{xt}}\nolimits}
\def\Sch{\mathit{Sch}}
\def\Mor{\mathop{\mathrm{Mor}}\nolimits}
\def\Ob{\mathop{\mathrm{Ob}}\nolimits}
\def\Sh{\mathop{\mathit{Sh}}\nolimits}
\def\NL{\mathop{N\!L}\nolimits}
\def\CH{\mathop{\mathrm{CH}}\nolimits}
\def\proetale{{pro\text{-}\acute{e}tale}}
\def\etale{{\acute{e}tale}}
\def\QCoh{\mathit{QCoh}}
\def\Ker{\mathop{\mathrm{Ker}}}
\def\Im{\mathop{\mathrm{Im}}}
\def\Coker{\mathop{\mathrm{Coker}}}
\def\Coim{\mathop{\mathrm{Coim}}}

% Boxtimes
%
\DeclareMathSymbol{\boxtimes}{\mathbin}{AMSa}{"02}

%
% Macros for moduli stacks/spaces
%
\def\QCohstack{\mathcal{QC}\!\mathit{oh}}
\def\Cohstack{\mathcal{C}\!\mathit{oh}}
\def\Spacesstack{\mathcal{S}\!\mathit{paces}}
\def\Quotfunctor{\mathrm{Quot}}
\def\Hilbfunctor{\mathrm{Hilb}}
\def\Curvesstack{\mathcal{C}\!\mathit{urves}}
\def\Polarizedstack{\mathcal{P}\!\mathit{olarized}}
\def\Complexesstack{\mathcal{C}\!\mathit{omplexes}}
% \Pic is the operator that assigns to X its picard group, usage \Pic(X)
% \Picardstack_{X/B} denotes the Picard stack of X over B
% \Picardfunctor_{X/B} denotes the Picard functor of X over B
\def\Pic{\mathop{\mathrm{Pic}}\nolimits}
\def\Picardstack{\mathcal{P}\!\mathit{ic}}
\def\Picardfunctor{\mathrm{Pic}}
\def\Deformationcategory{\mathcal{D}\!\mathit{ef}}
\title{A resolution computes the cohomology of the product total complex}
\author{The Stacks Project Authors}
\date{stacks-project:homology.tex:local-snapshot-20261012; GNU FDL 1.2 or later}
\begin{document}
\maketitle
\noindent\textit{Editorial note.} All delivered mathematical text is the author's own, copied from the pinned Stacks Project chapter file homology.tex; no statement or proof step was rewritten, reordered or removed.\par\medskip
\noindent\textit{Editorial note.} The proof is the author's complete argument; every step the author writes out is inside the counted span and no closing sentence is a placeholder.\par\medskip
\noindent\textit{Editorial note.} Every labelled reference occurring inside the counted spans is carried in assembly.external\_references as a non-load-bearing navigation pointer (same-chapter ones at their definition site, cross-chapter ones at the invocation inside the counted span), so no load-bearing mathematics is replaced by a stub. No bibliography citation occurs inside the counted spans.\par\medskip
\noindent\textit{Editorial note.} The same-name check for this card scanned every staged Stacks chapter of this repository -- all sources/stacks-*/pinned-original-source.tex files, the card's own chapter excluded -- and the labels found are recorded in dedup.same\_name\_cross\_chapter on the card.\par\medskip
\noindent\textit{Editorial note.} The preamble is the chapter's own native preamble excerpt already used for this source: the comment-type reference, slogan and history environments (native preamble lines 8-13), the Xy-pic setup (native preamble lines 14-18) and the native theorem declarations with the mathematical macros (native preamble lines 154-217, the actual last line of the file).\par\medskip
\noindent\textit{Allowed standard prerequisites (editorial).}\begin{itemize}
\item The product total complex of a double complex. Same chapter, the construction of the product total complex of a double complex, recalled in the proof. Supplies the total complex appearing in the statement and its differential.
\end{itemize}\medskip
\begin{lemma}
\label{lemma-resolution-gives-qis}
Let $M^\bullet$ be a complex of abelian groups. Let
$$
\ldots \to A_2^\bullet \to A_1^\bullet \to A_0^\bullet \to M^\bullet \to 0
$$
be an exact complex of complexes of abelian groups. Set $A^{p, q} = A_{-p}^q$
to obtain a double complex. Let $\text{Tot}_\pi(A^{\bullet, \bullet})$
be the product total complex associated to the double complex (see proof).
Then the map $\text{Tot}_\pi(A^{\bullet, \bullet}) \to M^\bullet$
induced by $A_0^\bullet \to M^\bullet$ is a quasi-isomorphism.
\end{lemma}

\begin{proof}
Abbreviating $T^\bullet = \text{Tot}_\pi(A^{\bullet, \bullet})$
we define
$$
T^n = \prod\nolimits_{p + q = n} A^{p, q} =
\prod\nolimits_{p + q = n} A_{-p}^q
\quad\text{with}\quad
\text{d}_{T^\bullet}^n =
\prod\nolimits_{n = p + q} (f_{-p}^q + (-1)^pd_{A_{-p}^\bullet}^q)
$$
where $f_p^\bullet : A_p^\bullet \to A_{p - 1}^\bullet$
are the maps of complexes in the lemma.
We will show that $T^\bullet$ is acyclic when
$M^\bullet$ is the zero complex. This will suffice by
the following trick. Set $B_n^\bullet = A_{n + 1}^\bullet$
and $B_0^\bullet = M^\bullet$. Then we have an exact sequence
$$
\ldots \to B_2^\bullet \to B_1^\bullet \to B_0^\bullet \to 0 \to 0
$$
as in the lemma. Let $S^\bullet = \text{Tot}_\pi(B^{\bullet, \bullet})$.
Then there is an obvious short exact sequence of complexes
$$
0 \to M^\bullet \to S^\bullet \to T^\bullet[1] \to 0
$$
and we conclude by the long exact cohomology sequence. Some details omitted.

\medskip\noindent
Assume $M^\bullet = 0$. We will show $H^0(T^\bullet) = 0$. The same argument
works for other degrees. Let $x =(x_n) \in \Ker(d_{T^\bullet})$
map to $\xi \in H^0(T^\bullet)$ with $x_n \in A^{-n, n} = A_n^n$.
Since $M^0 = 0$ we find that $x_0 = f_1^0(y_0)$ for some $y_0 \in A_1^0$.
Then $x_1 - d^0_{A_1^\bullet}(y_0) = f_2^1(y_1)$
because it is mapped to zero by $f_1^1$ as $x$ is a cocycle.
for some $y_1 \in A_2^1$. Continuing, using induction, we find
$y = (y_i) \in T^{-1}$ with $d_{T^\bullet}(y) = x$ as desired.
\end{proof}

\end{document}
